\documentclass[10pt,a4paper]{amsart}
\usepackage[margin=19mm]{geometry}
\usepackage{amsmath,amssymb,mathtools}
\usepackage[scr=rsfs,cal=euler]{mathalfa}
\usepackage{xfrac}
\usepackage[unicode,hidelinks,bookmarks=false]{hyperref}
\newcommand{\ie}{i.e.\ }
\newcommand{\cf}{cf.\ }
\newcommand{\resp}{respectively\ }
\newcommand{\Subsection}{\subsection}
\providecommand{\nobreakdash}{\nobreak\hbox{-}}
\setlength{\emergencystretch}{2em}
\multlinegap0pt
\usepackage{bm}
\usepackage{amssymb}
\usepackage{enumerate}
\usepackage{float}
\usepackage{mathtools}
\usepackage{xcolor}
\usepackage[normalem]{ulem}
\usepackage[shortlabels]{enumitem}
\newtheorem{lemma}{Lemma}
\newcommand{\A}{\mathbb{A}}
\newcommand{\Q}{\mathbb{Q}}
\newcommand{\edit}[1]{\textcolor{black}{#1}}
\newcommand{\ord}{\operatorname{ord}}
\newcommand{\val}{\vec{\boldsymbol{\alpha}}}
\newcommand{\veb}{\vec{\boldsymbol{\beta}}}
\title{Structure of (Fine) Mordell--Weil Groups}
\author{Rusiru Gambheera, Debanjana Kundu}
\date{Source: arXiv:2507.20341v3 (CC BY 4.0)}
\begin{document}
\maketitle

\noindent\emph{Source and licence.} Structure of (Fine) Mordell--Weil Groups. Version arXiv:2507.20341v3 (CC BY 4.0), distributed by arXiv under the Creative Commons Attribution 4.0 International licence, http://creativecommons.org/licenses/by/4.0/, as recorded in the arXiv licence metadata for this exact version and shown on the abstract page https://arxiv.org/abs/2507.20341v3. Full licence notice: \texttt{licenses/arxiv-2507.20341-CC-BY-4.0.txt}.\par\smallskip

\noindent\textit{Editorial scope.} Counted unit: the article's lemma 1 (source0.tex lines 755--765). The article's complete original proof of it is delivered with it (source0.tex lines 767--824, one \texttt{proof} environment of 58 lines). No mathematics is written, repaired or re-derived: the statement and the proof are the article's own lines, in the article's own notation, and no retained line is substituted or re-flowed.  No further source-local context is needed: every cross-reference of every retained line resolves inside the two delivered spans.  Nothing was omitted from any retained span, so the package carries no omission record.  No retained line contains a citation, so the wrapper carries no bibliography and no bibliography entry.  No retained line contains a figure environment, an image-inclusion command or drawing code, so no image is delivered and the package declares no asset dependency.  Editorial formatting only, in addition: an AMS \texttt{amsart} preamble that restates the article's own declarations and macros used by the retained lines, namely the environments \texttt{lemma} on separate counters, together with the standard packages those lines need.  The retained environments print in this document as Lemma 1 in order of appearance, whereas the article prints the same items with the numbering of its own full document.  The complete native source of the article as distributed in the arXiv e-print for this exact version is bundled unchanged as \texttt{sources/182278/source0.tex}.
\begin{lemma}
\label{lemma 1 for Prop 2 of note 1}
Suppose that $\val, \veb\in \A$.
Then
\[
W_{\veb}^{G_{\val}} = \begin{cases}
W_{\veb} & \text{ if } \veb \preccurlyeq \val\\
0 & \text{ otherwise}.
\end{cases}
\]
\end{lemma}

\begin{proof}
Consider the element
\[
\boldsymbol{w} = w_1 \otimes \ldots \otimes w_r \in W_{\veb},
\]
where each $w_i \in V_{i, \beta_i}$.
Also, let 
\[
g = (g_1 , \ldots, g_r)\in G_{\val}
\]
where each $g_i \in G^{{(i)}^{p_i^{\alpha_i}}}$.
Hence, each $g_i = \sigma_i^{k_{\edit{i}}p_i^{\alpha_i}}$ for some positive integer $k_{\edit{i}}$.
Via the isomorphism
\[
V_{i,\beta_i} \simeq \frac{\Q[x]}{\langle \omega(x^{p_i^{\beta_i - 1}}, p_i) \rangle}, 
\]
the action of $\sigma_i$ on $V_{i,\beta_i}$ corresponds to multiplication by $x$.
Thus, $\sigma_i^{k_{\edit{i}}p_i^{\alpha_i}}$ corresponds to $\left(x^{{p_i}^{\alpha_i}}\right)^{k_{\edit{i}}}$.

By definition, if $\veb\preccurlyeq \val$ then $\beta_i \leq \alpha_i$ for each $i$.
Thus
\[
\omega(x^{p_i^{\beta_i - 1}}, p_i) \mid x^{p_i^{\alpha_i}} - 1
\]
and
\[
\left( x^{p_i^{\alpha_i}} \right)^{k_{\edit{i}} }\equiv 1 \pmod{\omega(x^{p_i^{\beta_i - 1}}, p_i)}
\]
This means $g_i \cdot w_i = w_i$ for all $i$ and $g \cdot \boldsymbol{w} = \boldsymbol{w}$ as desired.

Next suppose that $\veb\not\preccurlyeq \val$ and that $\boldsymbol{w}\in W_{\veb}^{G_{\val}}$.
It is clear from the definition of partial ordering that there exists some $j$ such that $\beta_j > \alpha_j$.
Consider the element
\[
h = (1, 1, \ldots, 1, \sigma_j^{p_j^{\alpha_j}}, 1, \ldots, 1)\in G_{\val}.
\]
Note that $\ord(h) = p_j^{n_j - \alpha_j}$.
Write $\mathsf{N}(h)$ to denote the norm element, i.e.,
\[
\mathsf{N}(h) = 1 + h + h^2 + \ldots + h^{p_j^{n_j - \alpha_j}-1} \in \Q[G].
\]
Then, 
\begin{equation}
\label{eqn lemma 2.2}
\mathsf{N}(h)\cdot \boldsymbol{w} = p_j^{n_j - \alpha_j}\boldsymbol{w} \edit{.}
\end{equation}
On the other hand if we consider \edit{an arbitrary} element $\boldsymbol{w'}\in W_{\veb}$ then
\[
\mathsf{N}(h)\cdot \boldsymbol{w'} = w'_1 \otimes \ldots \otimes \sum_{k_{\edit{j}}=0}^{p_j^{n_j - \alpha_j}-1} \sigma_j^{k_{\edit{j}} p_j^{\alpha_j}}\cdot w'_j\otimes \ldots \otimes w'_r ,
\]
But, $\displaystyle\sum_{k_{\edit{j}}=0}^{p_j^{n_j - \alpha_j}-1} \sigma_j^{k_{\edit{j}} p_j^{\alpha_j}}$ acts on $V_{j,\beta_j} \simeq \frac{\Q[x]}{\langle \omega(x^{p_j^{\beta_j - 1}}, p_j) \rangle}$ as
\[
1 + x^{p_j^{\alpha_j}} + x^{2p_j^{\alpha_j}} + \ldots + x^{p_j^{n_j} - p_j^{\alpha_j}} = \frac{x^{p_j^{n_j}}-1}{x^{p_j^{\alpha_j}}-1}.
\]
But, $n_j \geq \beta_j > \alpha_j$, so $\omega(x^{p_j^{\beta_j - 1}}, p_j)$ divides $\frac{x^{p_j^{n_j}}-1}{x^{p_j^{\alpha_j}}-1}$.
Thus, $\mathsf{N}(h)\cdot \boldsymbol{w'}=0$.
By \eqref{eqn lemma 2.2}, it follows that $\mathsf{N}(h)\cdot \boldsymbol{w}\edit{=0}$ and hence $\boldsymbol{w}=0$.
\end{proof}

\end{document}
